
%% Use the "normalphoto" option if you want a normal photo instead of cropped to a circle
% \documentclass[10pt,a4paper,normalphoto]{altacv}

\documentclass[10pt,a4paper,ragged2e,withhyper]{altacv}
%% AltaCV uses the fontawesome5 and simpleicons packages.
%% See http://texdoc.net/pkg/fontawesome5 and http://texdoc.net/pkg/simpleicons for full list of symbols.

% Change the page layout if you need to
\geometry{left=1.25cm,right=1.25cm,top=1.5cm,bottom=1.5cm,columnsep=1.2cm}

% The paracol package lets you typeset columns of text in parallel
\usepackage{paracol}

\usepackage{ragged2e}

% xelatex -shell-escape -output-driver="xdvipdfmx -z 0" sample.tex
\ifxetexorluatex
  % If using xelatex or lualatex:
  \setmainfont{Roboto Slab}
  \setsansfont{Lato}
  \renewcommand{\familydefault}{\sfdefault}
\else
  % If using pdflatex:
  \usepackage[rm]{roboto}
  \usepackage[defaultsans]{lato}
  % \usepackage{sourcesanspro}
  \renewcommand{\familydefault}{\sfdefault}
\fi

% Change the colours if you want to
\definecolor{SlateGrey}{HTML}{2E2E2E}
\definecolor{LightGrey}{HTML}{666666}
\definecolor{DarkPastelRed}{HTML}{450808}
\definecolor{PastelRed}{HTML}{8F0D0D}
\definecolor{GoldenEarth}{HTML}{E7D192}
\colorlet{name}{black}
\colorlet{tagline}{PastelRed}
\colorlet{heading}{DarkPastelRed}
\colorlet{headingrule}{GoldenEarth}
\colorlet{subheading}{PastelRed}
\colorlet{accent}{PastelRed}
\colorlet{emphasis}{SlateGrey}
\colorlet{body}{LightGrey}

% Change some fonts, if necessary
\renewcommand{\namefont}{\Huge\rmfamily\bfseries}
\renewcommand{\personalinfofont}{\footnotesize}
\renewcommand{\cvsectionfont}{\LARGE\rmfamily\bfseries}
\renewcommand{\cvsubsectionfont}{\large\bfseries}


% Change the bullets for itemize and rating marker
% for \cvskill if you want to
\renewcommand{\cvItemMarker}{{\small\textbullet}}
\renewcommand{\cvRatingMarker}{\faCircle}
% ...and the markers for the date/location for \cvevent
% \renewcommand{\cvDateMarker}{\faCalendar*[regular]}
% \renewcommand{\cvLocationMarker}{\faMapMarker*}


% If your CV/résumé is in a language other than English,
% then you probably want to change these so that when you
% copy-paste from the PDF or run pdftotext, the location
% and date marker icons for \cvevent will paste as correct
% translations. For example Spanish:
% \renewcommand{\locationname}{Ubicación}
% \renewcommand{\datename}{Fecha}

\input{pubs-num.tex}

\addbibresource{sample.bib}

\begin{document}
\name{Lucas Noirot}

\tagline{Ingénieur R\&D en Algorithmique \& Mathématiques Appliquées}
%% You can add multiple photos on the left or right
\photoR{4cm}{img_7056}
%\photoL{2.8cm}{hovercode}

\personalinfo{%
  \email{lucas.noirot.contact@gmail.com}
  %\phone{(+33) 685263667}
  %\mailaddress{Åddrésş, Street, 00000 Cóuntry}
  \location{Montpellier, France}
  \homepage{crowreed.github.io}
  \linkedin{lucas-petit-dev}
  \github{Crowreed}
  %\orcid{0000-0000-0000-0000}
  %% You can add your own arbitrary detail with
  %% \printinfo{symbol}{detail}[optional hyperlink prefix]
  % \printinfo{\faPaw}{Hey ho!}[https://example.com/]

  %% Or you can declare your own field with
  %% \NewInfoFiled{fieldname}{symbol}[optional hyperlink prefix] and use it:
  % \NewInfoField{gitlab}{\faGitlab}[https://gitlab.com/]
  % \gitlab{your_id}
  %%
  %% For services and platforms like Mastodon where there isn't a
  %% straightforward relation between the user ID/nickname and the hyperlink,
  %% you can use \printinfo directly e.g.
  % \printinfo{\faMastodon}{@username@instace}[https://instance.url/@username]
  %% But if you absolutely want to create new dedicated info fields for
  %% such platforms, then use \NewInfoField* with a star:
  % \NewInfoField*{mastodon}{\faMastodon}
  %% then you can use \mastodon, with TWO arguments where the 2nd argument is
  %% the full hyperlink.
  % \mastodon{@username@instance}{https://instance.url/@username}
}

%% Depending on your tastes, you may want to make fonts of itemize environments slightly smaller
% \AtBeginEnvironment{itemize}{\small}

%% Set the left/right column width ratio to 6:4.
\columnratio{0.3}


\begin{paracol}{2}



%%%%%%%%%%%%%%%%%%%%%%%%%%%%%%%%%%%%%%%%%%%%%%%%%%%%%%%%%%%%%%%%%
\makecvheaderright

\cvsection{À propos de moi}
\vspace{-1em}
\begin{quote}
    \begin{justify}
        Passionné par l'informatique et les mathématiques, je trouve mon équilibre dans l'abstraction, où les idées prennent forme et deviennent concrètes.
    \end{justify}
\end{quote}



\cvsection{Compétences}



\cvsubsection{Rédaction}

\cvtag{\LaTeX, MarkDown} \hfill\cvtag{Pédagogie}

\cvsubsection{Stack Technique}

\cvtag{C/C++}\hfill
\cvtag{HTML5/CSS3/JavaScript}\hfill
\cvtag{Java}\hfill
\cvtag{Python}\hfill
\cvtag{SageMath}

\vspace{1em}
\cvsubsection{Compétences informatiques}

\cvtag{Linux, Windows} \hfill
\cvtag{Git}


%%%%%%%%%%%%%%%%%%%%%%%%%%%%%%%%%%%%%%%%%%%%%%%%%%%%%%%%%%%%%%%%%

\cvsection{Langues}

\textbf{Français} -- natif

\vspace{1em}

\textbf{Anglais} -- Niveau C1 · Maîtrise à l’écrit et à l’oral · Anglais technique courant


\cvsection{Autre}


\cvsubsection{Qualités personnelles}


\cvtag{Rigoureux} \hfill
\cvtag{Autodidacte}\\
\cvtag{Pro-actif} \hfill
\cvtag{Abstraction}

\vspace{1em}
\cvsubsection{Centres d'intérêt}

\cvtag{Origami} \hfill
\cvtag{Randonnée}\\
\cvtag{Astronomie} \hfill
\cvtag{Jeu vidéo}
\medskip

%%%%%%%%%%%%%%%%%%%%%%%%%%%%%%%%%%%%%%%%%%%%%%%%%%%%%%%%%%%%%%%%%%%%%%%%%%%%%%%%%%%%%%%%%%%%%%%%%%%%%%%%%%%%%%%%%%%%%%%%%%%%%%%%%%
%%%%%%%%%%%%%%%%%%%%%%%%%%%%%%%%%%%%%%%%%%%%%%%%%%%%%%%%%%%%%%%%%%%%%%%%%%%%%%%%%%%%%%%%%%%%%%%%%%%%%%%%%%%%%%%%%%%%%%%%%%%%%%%%%%
%%%%%%%%%%%%%%%%%%%%%%%%%%%%%%%%%%%%%%%%%%%%%%%%%%%%%%%%%%%%%%%%%%%%%%%%%%%%%%%%%%%%%%%%%%%%%%%%%%%%%%%%%%%%%%%%%%%%%%%%%%%%%%%%%%
%%%%%%%%%%%%%%%%%%%%%%%%%%%%%%%%%%%%%%%%%%%%%%%%%%%%%%%%%%%%%%%%%%%%%%%%%%%%%%%%%%%%%%%%%%%%%%%%%%%%%%%%%%%%%%%%%%%%%%%%%%%%%%%%%%

\switchcolumn

\makecvheaderleft


\cvsection{Formations}

\cveventt{Master en Informatique - Algorithmique}
         {Université de Montpellier}
         {2023 – 2025}
         {Complexité · Cryptographie · Calcul formel · Optimisation}

\vspace{0.7em}

\cveventt{Master 1 en Mathématiques}
         {Université de Bourgogne}
         {2022 – 2023}
         {Mention Assez bien}

\vspace{0.7em}

\cveventt{Licence en Mathématiques}
         {Université de Bourgogne}
         {2019 – 2022}
         {Classé 5\textsuperscript{e} sur 55}

\vspace{0.9em}

\cveventt{Baccalauréat scientifique}
         {Lycée Denis Diderot — Langres}
         {2016 – 2019}
         {Mention TB · Section Euro} % %Mention Très Bien · Section européenne · Spécialité sciences de l’ingénieur

%%%%%%%%%%%%%%%%%%%%%%%%%%%%%%%%%%%%%%%%%%%%%%%%%%%%%%%%%%%%%%%%%

\cvsection{Expérience}

\cveventt{Stage de recherche en cryptographie (4 mois)}{Laboratoire d’Informatique, de Robotique et de Microélectronique de Montpellier  – Équipe ECO}
         {2025}
         {R\&D algorithmique sous SageMath . Présentation technique en anglais (1h30)}

\vspace{0.9em}

\cveventt{Professeur particulier}
         {Dijon, freelance}
         {2022}
         {Mathématiques et physique niveau collègue et lycée}

\vspace{0.9em}

\cveventt{Opérateur de production}{Langres, Freudenberg}{Juillet 2021}{Travail d'été en équipe en environnement industriel}


\cvsection{Projets personnels}

\cvevent{Chat chiffré en C++ avec cryptographie post-quantique}{Conception et implémentation d’un chat chiffré de bout en bout, intégrant un protocole de cryptographie post-quantique fondé sur les réseaux euclidiens.}{}{}
\vspace{0.9em}

\cvevent{Articles de vulgarisation scientifique}{Rédaction d’articles pédagogiques expliquant les principes de la cryptographie post-quantique et des réseaux euclidiens.}{}{}
\vspace{0.9em}

\cvevent{Site web personnel (HTML / CSS / JavaScript)}{Développement d’une interface de présentation de projets scientifiques et techniques.}{}{}

% use ONLY \newpage if you want to force a page break for
% ONLY the current column
\newpage
%% Switch to the right column. This will now automatically move to the second
%% page if the content is too long.

\end{paracol}

\end{document}